\documentclass[11pt]{article}
\usepackage[a4paper,margin=1in]{geometry}
\usepackage{amsmath,amssymb,amsthm,mathtools,bm}
\usepackage{physics}
\usepackage{hyperref}

\newcommand{\E}{\mathbb{E}}
\DeclareMathOperator{\Cov}{Cov}
\DeclareMathOperator{\Var}{Var}

\title{Gaussian $J$-Integral in p=1 QAOA for SK and the Spin-Glass Overlap Order Parameter}
\author{}
\date{}

\begin{document}
	\maketitle
	
	\section*{1. Setup and notation}
	We consider the SK cost (your code normalization):
	\begin{equation}
		C_J(z)
		=
		\frac{1}{\sqrt{N}}
		\sum_{j<k} J_{jk}\,\sigma_j(z)\sigma_k(z),
		\qquad N:=n+1,
		\label{eq:CJ}
	\end{equation}
	where $\sigma_j(z)\in\{\pm 1\}$ is the spin associated with bitstring/basis state $z$.
	The couplings are i.i.d.
	\begin{equation}
		J_{jk}\sim\mathcal N(0,1),\qquad j<k,
		\label{eq:J-gaussian}
	\end{equation}
	and independent across edges.
	
	After previous algebraic steps, one integral appearing term-by-term in the $x,y,z$ sum is
	\begin{equation}
		I_{xyz}
		:=
		\int U_C^\dagger(x;J)\,\frac{C_J(z)}{N}\,U_C(y;J)\,d\mu(J),
		\label{eq:Ixyz-def}
	\end{equation}
	with
	\begin{equation}
		U_C(y;J)=e^{-i\gamma C_J(y)},\qquad
		U_C^\dagger(x;J)=e^{+i\gamma C_J(x)}.
	\end{equation}
	Hence
	\begin{equation}
		I_{xyz}
		=
		\frac{1}{N}\,\E\!\left[C_z\,e^{i\gamma(C_x-C_y)}\right],
		\label{eq:Ixyz-mgf-form}
	\end{equation}
	where for brevity $C_u:=C_J(u)$.
	
	\section*{2. Gaussian evaluation identity}
	Define
	\begin{equation}
		X:=C_z,\qquad Y:=C_x-C_y.
	\end{equation}
	Since all $C_u$ are linear forms in Gaussian $J$, the vector $(X,Y)$ is jointly Gaussian and centered.
	A standard joint-Gaussian identity gives
	\begin{equation}
		\E\!\left[X e^{itY}\right]
		=
		(it)\,\Cov(X,Y)\,\exp\!\left(-\frac{t^2}{2}\Var(Y)\right).
		\label{eq:gaussian-identity}
	\end{equation}
	Applying $t=\gamma$ yields
	\begin{equation}
		I_{xyz}
		=
		\frac{i\gamma}{N}\,\Cov(C_z, C_x-C_y)
		\exp\!\left(-\frac{\gamma^2}{2}\Var(C_x-C_y)\right).
		\label{eq:Ixyz-cov-var}
	\end{equation}
	This is already the closed $J$-integral in covariance/variance form.
	
	\section*{3. Compute covariances from SK normalization}
	Let
	\begin{equation}
		a_{jk}(u):=\sigma_j(u)\sigma_k(u),\qquad
		C_u=\frac{1}{\sqrt N}\sum_{j<k}J_{jk}a_{jk}(u).
	\end{equation}
	Using independence and $\E[J_{jk}J_{\ell m}]=\delta_{(jk),(\ell m)}$,
	\begin{equation}
		\Cov(C_u,C_v)
		=
		\frac{1}{N}\sum_{j<k} a_{jk}(u)a_{jk}(v).
		\label{eq:cov-edge-sum}
	\end{equation}
	Introduce the overlap
	\begin{equation}
		q_{uv}:=\frac{1}{N}\sum_{j=1}^{N}\sigma_j(u)\sigma_j(v)\in[-1,1].
		\label{eq:q-def}
	\end{equation}
	Using
	\begin{equation}
		\sum_{j<k} b_j b_k = \frac12\!\left[\left(\sum_j b_j\right)^2-\sum_j b_j^2\right],
		\quad b_j:=\sigma_j(u)\sigma_j(v),\quad b_j^2=1,
	\end{equation}
	we get
	\begin{equation}
		\Cov(C_u,C_v)
		=
		\frac{N q_{uv}^2-1}{2}.
		\label{eq:cov-q}
	\end{equation}
	Therefore
	\begin{align}
		\Cov(C_z, C_x-C_y)
		&=\Cov(C_z,C_x)-\Cov(C_z,C_y)
		\\
		&=\frac{N}{2}\left(q_{zx}^2-q_{zy}^2\right),
		\label{eq:cov-diff-q}
	\end{align}
	(the $-\tfrac12$ cancels).
	Also
	\begin{align}
		\Var(C_x-C_y)
		&=\Var(C_x)+\Var(C_y)-2\Cov(C_x,C_y)
		\\
		&=\frac{N-1}{2}+\frac{N-1}{2}-2\cdot\frac{N q_{xy}^2-1}{2}
		\\
		&=N\left(1-q_{xy}^2\right).
		\label{eq:var-diff-q}
	\end{align}
	Plugging \eqref{eq:cov-diff-q} and \eqref{eq:var-diff-q} into \eqref{eq:Ixyz-cov-var}:
	\begin{equation}
		\boxed{
			I_{xyz}
			=
			\frac{i\gamma}{2}\left(q_{zx}^2-q_{zy}^2\right)
			\exp\!\left[-\frac{\gamma^2}{2}N\left(1-q_{xy}^2\right)\right]
		}
		\qquad (N=n+1).
		\label{eq:Ixyz-final}
	\end{equation}
	This is the explicit Gaussian $J$-integral result for each $(x,y,z)$ term.
	
	\section*{4. Where this sits in your current formula}
	With kernel factors already pulled outside the integral, each summand is
	\begin{equation}
		\overline{K_\beta(z,x)}\,K_\beta(z,y)\,I_{xyz},
	\end{equation}
	so the expectation becomes
	\begin{equation}
		\E_J[\langle\cdots\rangle]
		=
		\left(\frac12\right)^N
		\sum_{x,y,z}
		\overline{K_\beta(z,x)}\,K_\beta(z,y)
		\cdot I_{xyz},
	\end{equation}
	with $I_{xyz}$ given by \eqref{eq:Ixyz-final}.
	This is exactly the point before doing the remaining combinatorial/mixer-kernel summation.
	
	\section*{5. What is the overlap $q$?}
	The overlap
	\begin{equation}
		q_{uv}=\frac1N\sum_j \sigma_j(u)\sigma_j(v)
	\end{equation}
	measures similarity between spin configurations $u,v$.
	Equivalent Hamming-distance form:
	\begin{equation}
		q_{uv}=1-\frac{2\,d_H(u,v)}{N}.
	\end{equation}
	So $q=1$ means identical states, $q=-1$ means fully opposite spins, and intermediate values quantify partial alignment.
	
	\section*{6. Spin-glass meaning (SK theory)}
	In spin-glass theory, $q$ is the fundamental order parameter:
	\begin{itemize}
		\item It labels similarity between different Gibbs samples (``replicas'').
		\item The distribution $P(q)$ is the object that distinguishes phases.
		\item Replica-symmetric regimes have a simpler overlap structure; full RSB has a nontrivial broad $P(q)$ and ultrametric organization of pure states.
	\end{itemize}
	In this QAOA/SK calculation, overlap appears directly because Gaussian cost correlations satisfy
	\begin{equation}
		\Cov(C_u,C_v)\propto q_{uv}^2,
	\end{equation}
	so all $J$-averaged phase factors and cost insertions reduce to overlap-dependent expressions.
	
	\section*{7. Practical interpretation for your proof pipeline}
	At your current stage:
	\begin{enumerate}
		\item Sum/integral order has been swapped (with integrability assumptions).
		\item Kernel factors are outside the $J$-integral.
		\item The remaining $J$-integral is now explicit via \eqref{eq:Ixyz-final}.
	\end{enumerate}
	The next proof task is therefore purely algebraic/combinatorial in $x,y,z$ plus trigonometric simplification through the mixer kernel structure.
	
\end{document}
